% §3.3 The Composition Theorem
%
% Cross-reference resolution status:
%
% EXPECTED UNRESOLVED (forward refs that resolve in later Tasks/Plans):
%   \Cref{sec:miller-rabin}       -- chapter 3 §3.4 (resolves in Plan 3 Task 5)
%   \Cref{sec:three-faces}        -- chapter 1 §1.1 (resolves in Plan 4)
%
% ALREADY RESOLVED (chapters 2 and 4, drafted in Plans 1 and 2;
%                   and §3.1-§3.2, drafted in Tasks 2-3 of Plan 3):
%   \Cref{sec:channel-composition}  -- chapter 2 §2.3
%   \Cref{eq:bsc-composition}       -- chapter 2 §2.3 BSC composition formula
%   \Cref{def:channel-composition}  -- chapter 2 §2.3 serial composition definition
%   \Cref{sec:and-under-noise}      -- chapter 3 §3.1
%   \Cref{eq:and-fnr}               -- §3.1 AND false-negative rate
%   \Cref{sec:or-not}               -- chapter 3 §3.2
%   \Cref{eq:or-fpr-direct}         -- §3.2 OR false-positive rate
%   \Cref{sec:kronecker}            -- chapter 4 §4.4
%   \Cref{thm:kronecker}            -- chapter 4 §4.4 Kronecker theorem
%   \Cref{sec:axiom-1}              -- chapter 4 §4.2
%   \Cref{ax:element-indep}         -- Axiom 1 formal statement
%   \Cref{def:bernoulli-bool}       -- chapter 4 §4.5
%
% NOTE ON THE BSC RECONCILIATION (plan Task 4 Step 5):
%   The composition theorem's formula eta_total = 1 - prod(1 - eta_i) gives
%   the probability of at least one error in the chain. For identity-with-error
%   composition (any single error corrupts the end-to-end result), this is the
%   correct formula. For BSC chains, an even number of flips cancels out, giving
%   the chapter 2 §2.3 formula 2*eps - 2*eps^2 (odd-flip probability) rather
%   than 2*eps - eps^2 (at-least-one-flip probability). The section makes this
%   distinction explicit. The composition theorem applies cleanly to one-sided
%   error channels and to the AND/OR gates; BSC composition requires the
%   §2.3 matrix-product formula.
%
% Build warnings for the two unresolved labels above are expected at this stage.
% They resolve when Plan 3 Task 5 and Plan 4 land.

\section{The Composition Theorem}
\label{sec:composition-theorem}

\Cref{sec:and-under-noise} showed that the AND of $k$ independent
$\Bernoulli{\BoolType}$ witnesses drives the false-positive rate down:
requiring all $k$ to agree on $\top$ makes a false $\top$ output unlikely.
\Cref{sec:or-not} showed the dual: the OR of $k$ witnesses drives the
false-negative rate down. Both sections closed with a preview of the
same structural pattern: the probability that the chain produces \emph{no}
error factors as a product.

This section makes that pattern precise. We state and prove the
\emph{composition theorem}: for $k$ independent identity-with-error stages
chained in series, the end-to-end error rate is
\[
  \eta_{\mathrm{total}} = 1 - \prod_{i=1}^{k}(1 - \eta_i),
\]
where $\eta_i$ is the per-stage error rate. We then examine what this formula
does and does not capture, paying particular attention to the BSC running
thread---because the BSC has a subtlety, a cancellation structure, that the
composition theorem does not capture. Understanding that gap is essential to
using the theorem correctly.

\subsection*{Setup: chaining identity-with-error stages}

The model is simple. We have an original Boolean value $b \in \{\top, \bot\}$.
A single \emph{identity-with-error} stage passes $b$ through a channel that
reports the correct value with probability $1 - \eta_i$ and reports the wrong
value with probability $\eta_i$. Call the stage's output $\hat{b}_i$. The stage
is not a BSC in general: we will say more shortly about what the error rate
$\eta_i$ means. For now, think of it as the probability that the stage makes
\emph{any} error, regardless of direction.

We chain $k$ such stages: $b \to \hat{b}_1 \to \hat{b}_2 \to \cdots \to
\hat{b}_k$. The final output $\hat{b}_k$ is what we observe. The question is:
what is the probability that $\hat{b}_k \neq b$?

For this question to have a clean answer, we need a clean independence
structure. We assume:

\begin{enumerate}
\item \textbf{Chain-stage independence.} The error event at stage $i$ is
  independent of the error event at every other stage $j \neq i$. This is
  a chain-level assumption: it says that the $k$ stages do not share
  error-inducing correlations along the pipeline.

\item \textbf{Each error is one-directional.} Once stage $i$ makes an error,
  the subsequent stages in the chain receive the wrong value and propagate it.
  Errors do not cancel. This is the identity-with-error regime: any single
  error in the chain corrupts the end-to-end output.
\end{enumerate}

The second assumption is worth flagging before the theorem, not after. It is
the source of the BSC subtlety we will return to.

\subsection*{The composition theorem}

\begin{theorem}[Composition theorem, identity-with-error chains]
\label{thm:composition}
Let $b$ be a Boolean value. Let $\hat{b}_1, \hat{b}_2, \ldots, \hat{b}_k$
be the successive outputs of $k$ chained identity-with-error stages, where
stage $i$ has error rate $\eta_i \in [0,1]$ and the error events across
stages are mutually independent.

Under the identity-with-error assumption (any single error in the chain
corrupts the end-to-end output), the end-to-end error rate satisfies
\begin{equation}
\label{eq:composition-theorem}
  \eta_{\mathrm{total}}
  \;=\; 1 - \prod_{i=1}^{k}(1 - \eta_i).
\end{equation}

Equivalently, letting $\tnrate_i = 1 - \eta_i$ denote the per-stage
\emph{correctness probability}, the end-to-end correctness probability is
\begin{equation}
\label{eq:composition-correctness}
  \tnrate_{\mathrm{total}}
  \;=\; \prod_{i=1}^{k} \tnrate_i.
\end{equation}
\end{theorem}

\begin{proof}
We prove the correctness-probability form~\eqref{eq:composition-correctness}
by induction on $k$, then recover~\eqref{eq:composition-theorem} by
complementation.

\medskip
\noindent\textbf{Base case, $k = 1$.} A single stage with error rate $\eta_1$
has correctness probability $\tnrate_1 = 1 - \eta_1$. The formula gives
$\prod_{i=1}^{1} \tnrate_i = \tnrate_1$. Correct.

\medskip
\noindent\textbf{Inductive step.} Assume the result holds for $k-1$ stages:
the first $k-1$ stages in the chain produce a correct output with probability
$\tnrate_1 \cdots \tnrate_{k-1}$. Now append stage $k$ with error rate
$\eta_k$.

Under the identity-with-error assumption, the final output $\hat{b}_k$ equals
$b$ if and only if \emph{both} of the following hold:
\begin{itemize}
  \item The output of the first $k-1$ stages is correct, with
    probability $\tnrate_1 \cdots \tnrate_{k-1}$ (by the inductive hypothesis).
  \item Stage $k$ does not make an error, with probability $\tnrate_k$.
\end{itemize}
By chain-stage independence, these two events are independent. Their joint
probability is:
\[
  \tnrate_{\mathrm{total}}
  \;=\; (\tnrate_1 \cdots \tnrate_{k-1}) \cdot \tnrate_k
  \;=\; \prod_{i=1}^{k} \tnrate_i.
\]

Taking the complement, the end-to-end error rate is
\[
  \eta_{\mathrm{total}} = 1 - \tnrate_{\mathrm{total}} = 1 - \prod_{i=1}^{k}(1 - \eta_i).
\qedhere
\]
\end{proof}

\noindent\textbf{Remark on the identity-with-error assumption.}
The inductive step rests on the phrase ``if and only if both.'' This is only
valid when any error at any stage is fatal to correctness. If stage $k$
receives an already-wrong value and makes an error, it \emph{could} accidentally
correct the previous error---returning to the right answer. The composition
theorem excludes this possibility by assumption. For channels where
cancellation is possible (notably the BSC), the formula~\eqref{eq:composition-theorem}
does not give the exact end-to-end error rate, and the matrix-product
computation of \Cref{sec:channel-composition} must be used instead. We return
to this point with numerical evidence in the BSC worked example below.

\subsection*{Connection to the AND and OR formulas}

Before the BSC example, let us confirm that the composition theorem unifies
the formulas already derived in \Cref{sec:and-under-noise,sec:or-not}.

\medskip
\noindent\textbf{AND as composition.}
In \Cref{sec:and-under-noise}, the false-negative rate of AND was:
\[
  \fnrate_\wedge
  \;=\; 1 - (1 - \fnrate_1)(1 - \fnrate_2).
\]
This is exactly~\eqref{eq:composition-theorem} with $k = 2$ and $\eta_i = \fnrate_i$.
Interpreting AND as a two-stage identity-with-error composition in the positive block:
stage 1 fails to deliver a $\top$ output with probability $\fnrate_1$; stage 2
independently fails with probability $\fnrate_2$; the AND of the two outputs is
$\top$ only if both succeed. The identity-with-error assumption is satisfied in
the positive block because if either stage makes a false negative (reports $\bot$
when the true value is $\top$), the AND output is $\bot$ regardless of the other
stage. There is no cancellation in the $\top$ direction.

Extending to $k$ independent $\Bernoulli{\BoolType}$ instances all having true
value $\top$, with false-negative rates $\fnrate_1, \ldots, \fnrate_k$: the
$k$-fold AND has false-negative rate
\[
  \fnrate_{\wedge,k} \;=\; 1 - \prod_{i=1}^{k}(1 - \fnrate_i).
\]
When all stages have the same rate $\fnrate$, this simplifies to
$1 - (1 - \fnrate)^k$.

\medskip
\noindent\textbf{OR as composition.}
In \Cref{sec:or-not}, the false-positive rate of OR was:
\[
  \fprate_\vee
  \;=\; 1 - (1 - \fprate_1)(1 - \fprate_2).
\]
This is~\eqref{eq:composition-theorem} with $k = 2$ and $\eta_i = \fprate_i$.
Interpreting OR as a two-stage identity-with-error composition in the negative
block: stage 1 makes a false positive with probability $\fprate_1$; stage 2
independently with probability $\fprate_2$; the OR output is $\bot$ only if
both avoid false positives. The identity-with-error assumption holds in the
$\bot$ direction because once either stage makes a false positive (reports
$\top$ when the true value is $\bot$), the OR output is $\top$ regardless.

Extending to $k$ stages with false-positive rates $\fprate_1, \ldots, \fprate_k$:
the $k$-fold OR has false-positive rate
\[
  \fprate_{\vee,k} \;=\; 1 - \prod_{i=1}^{k}(1 - \fprate_i).
\]

\subsection*{The BSC running thread: where the formula applies and where it does not}

The binary symmetric channel (BSC) from \Cref{sec:channel-composition} is the
principal running example for composition in this book. We use it here to
illustrate both what the composition theorem captures and, critically, where
its identity-with-error assumption breaks down.

\medskip
\noindent\textbf{What chapter 2 §2.3 derived.}
\Cref{sec:channel-composition} computed the crossover probability of two
BSCs in series using the matrix product
$Q_{\mathrm{total}} = Q_2 Q_1$. The result (\Cref{eq:bsc-composition}) was:
\[
  \fprate_{\mathrm{total}}
  \;=\; \fprate_1 + \fprate_2 - 2\fprate_1\fprate_2.
\]
For equal crossovers $\fprate_1 = \fprate_2 = \fprate$, this becomes:
\begin{equation}
\label{eq:bsc-two-stage-correct}
  \fprate_{\mathrm{total}}
  \;=\; 2\fprate - 2\fprate^2.
\end{equation}

\medskip
\noindent\textbf{What the composition theorem gives.}
If we naively apply~\eqref{eq:composition-theorem} with $\eta_i = \fprate$ for
both stages and $k = 2$, we obtain:
\begin{equation}
\label{eq:bsc-two-stage-naive}
  \eta_{\mathrm{total}}
  \;=\; 1 - (1 - \fprate)^2
  \;=\; 2\fprate - \fprate^2.
\end{equation}

\medskip
\noindent\textbf{These two formulas are different.}
\[
  \underbrace{2\fprate - 2\fprate^2}_{\text{correct BSC composition, §2.3}}
  \;\neq\;
  \underbrace{2\fprate - \fprate^2}_{\text{naive §3.3 formula}}.
\]
For $\fprate = 0.1$:
\begin{align*}
  \text{correct (§2.3): } & 2(0.1) - 2(0.01) = 0.18, \\
  \text{naive (§3.3): }   & 2(0.1) - (0.01) = 0.19.
\end{align*}
The naive formula overestimates the BSC's end-to-end error rate.

The discrepancy is not a bug in either formula. It reveals the identity-with-error
assumption in action.

\medskip
\noindent\textbf{The cancellation structure of the BSC.}
In a BSC, a ``flip'' swaps $\top \leftrightarrow \bot$. Two flips cancel:
$\top \xrightarrow{\text{flip}} \bot \xrightarrow{\text{flip}} \top$. The
end-to-end output is wrong if and only if an \emph{odd} number of flips
occurred across the chain. The end-to-end output is right if an \emph{even}
number of flips occurred (including zero).

For $k = 2$ stages with equal crossover $\fprate$, the parity decomposition is:
\begin{align*}
  \Pr[\text{exactly 0 flips}] &= (1-\fprate)^2, \\
  \Pr[\text{exactly 1 flip}]  &= 2\fprate(1-\fprate), \\
  \Pr[\text{exactly 2 flips}] &= \fprate^2.
\end{align*}
The end-to-end error probability is the odd-flip probability:
\[
  \Pr[\text{odd number of flips}]
  = 2\fprate(1-\fprate)
  = 2\fprate - 2\fprate^2.
\]
This matches the §2.3 formula~\eqref{eq:bsc-two-stage-correct} exactly. The
``at-least-one-flip'' probability that the naive §3.3 formula computes is
\[
  \Pr[\text{at least one flip}]
  = 1 - (1-\fprate)^2
  = 2\fprate - \fprate^2,
\]
which includes the event that both stages flip---an event that \emph{cancels}
to correctness. The composition theorem counts this event as an error; the BSC
does not. Hence the overestimate.

\medskip
\noindent\textbf{Summary of the distinction.}
The composition theorem formula $\eta_{\mathrm{total}} = 1 - \prod(1 - \eta_i)$
answers the question: \emph{what is the probability that at least one stage in
the chain makes an error?} For the identity-with-error model---where a single
error corrupts the end-to-end output irrecoverably---this is also the
end-to-end error rate. But for the BSC, two errors cancel, so the relevant
question is: \emph{what is the probability that an odd number of stages flip?}
The BSC composition formula from \Cref{sec:channel-composition} answers this
second question.

The two questions agree only in special cases:
\begin{itemize}
  \item \textbf{One-sided error channels: FNR $= 0$.} If $\fnrate = 0$ (a
    positive approximate set), then stage $i$ can only make errors in one
    direction: a false positive cannot be ``uncorrected'' by a later false
    positive, because the false positive takes a $\bot$ input to $\top$, and
    a subsequent channel with FNR $= 0$ will report $\top$ correctly. There
    is no cancellation in the positive direction. The composition theorem
    formula gives the exact accumulated false-positive rate.

  \item \textbf{One-sided error channels: FPR $= 0$.} Symmetrically, if
    $\fprate = 0$, the only errors are false negatives, and two false negatives
    do not cancel (a missed positive remains missed). The composition theorem
    applies.

  \item \textbf{AND / OR gates.} As shown above, AND with FNR rates and OR
    with FPR rates are both in the identity-with-error regime in their
    respective dominant error direction. No cancellation is possible.

  \item \textbf{Miller-Rabin's $k$-round test.} The Miller-Rabin witness test
    has FNR $= 0$ per round (it never falsely declares a prime composite).
    The only error is a false positive, and a false positive cannot be
    ``uncanceled'' by subsequent rounds---each round independently either
    detects the compositeness or doesn't. The composition theorem applies
    directly: $\fprate_k \leq (1/4)^k$. We work through this in
    \Cref{sec:miller-rabin}.
\end{itemize}

For the BSC, neither condition holds: the BSC has $\fprate = \fnrate$, and two
flips cancel. The correct composition formula is the one in
\Cref{sec:channel-composition}: matrix product, not the composition theorem
of this section. This is not a contradiction; it is two different theorems
answering two different questions about two different models.

\subsection*{The channel-matrix product view}

The matrix-product computation in \Cref{sec:channel-composition} gives the
BSC composition formula (with cancellation): two flips that cancel return the
chain to the correct answer, and the resulting end-to-end crossover is
$\fprate_1 + \fprate_2 - 2\fprate_1\fprate_2$, not $1 - (1-\fprate_1)(1-\fprate_2)$.
The composition theorem of this section applies to a different underlying
model: each stage either preserves correctness or destroys it, with no
possibility of recovery (``once lost, stays lost''). The BSC's bidirectional
flip structure is incompatible with this monotone correctness assumption.
The two formulas describe different physical processes; neither is wrong,
and they answer different questions about different models.

\subsection*{Sister structure: Kronecker for parallel queries}

The composition theorem governs \emph{serial} chains: one query passes through
$k$ stages in sequence. The Kronecker factorization theorem
(\Cref{thm:kronecker}) governs \emph{parallel} queries: $m$ queries on $m$
distinct elements, each through its own copy of the channel.

In the serial case (this section), correctness probabilities multiply along
the chain:
\[
  \tnrate_{\mathrm{total}} \;=\; \tnrate_1 \cdot \tnrate_2 \cdots \tnrate_k.
\]

In the parallel case (\Cref{sec:kronecker}), the joint confusion matrix factors
as a Kronecker product:
\[
  Q^{(m)} \;=\; Q_1 \otimes Q_2 \otimes \cdots \otimes Q_m.
\]

These two factorizations arise from different independence structures. The
serial factorization uses \emph{chain-stage independence}: the error events at
distinct stages of the pipeline are independent. The Kronecker factorization
uses \emph{Axiom 1} (\Cref{ax:element-indep}): the error events at distinct
input elements are independent. The two independence assumptions can both hold
simultaneously---and in well-designed systems, they do---but they are logically
distinct.

An important practical consequence: when $k$ parallel pipelines each consist of
multiple stages in series, the combined channel matrix is
\[
  \left(Q_2^{(1)} Q_1^{(1)}\right)
  \otimes
  \left(Q_2^{(2)} Q_1^{(2)}\right)
  \otimes \cdots,
\]
as \Cref{sec:channel-composition} noted. The serial composition (matrix product)
applies \emph{within} each pipeline; the Kronecker product applies \emph{across}
parallel pipelines. The composition theorem of this section computes the
correctness probability of each pipeline's serial chain; the Kronecker theorem
assembles the joint confusion matrix for the $m$ parallel pipelines together.

\subsection*{A note on chain-stage independence vs.\ element-wise independence}

The composition theorem assumes chain-stage independence. Axiom~1
(\Cref{ax:element-indep}) guarantees element-wise independence: error events
at distinct input elements are independent. These are closely related but
operationally distinct.

In the Miller-Rabin example of \Cref{sec:miller-rabin}, the $k$ witness
evaluations are drawn from a pool of witnesses. If we treat each witness as a
distinct element of a predicate over the witness pool---which is the right
recasting, as \Cref{sec:miller-rabin} details---then Axiom 1 licenses
independence across witnesses. That independence is also chain-stage
independence for the $k$-AND composition: the $k$ AND stages each involve a
distinct element, and their error events are independent by Axiom 1. In this
setting, Axiom 1 implies chain-stage independence, and the composition theorem
follows.

In contrast, consider a pipeline where stage $i$ and stage $j$ query the
\emph{same} element of the same Bernoulli predicate. If the implementation is
deterministic (same input, same output), the two stages are perfectly
correlated, not independent. Axiom 1 does not apply across repeated queries of
the same element (Chapter 16 discusses this negative space). The composition
theorem would not apply either: the chain-stage independence assumption fails
when stages share elements.

The honest position is this: chain-stage independence is a structural property
of the pipeline, not a consequence of Axiom 1 alone. In well-behaved
implementations where each stage of the pipeline queries a distinct element,
Axiom 1 implies chain-stage independence and the composition theorem holds.
In pipelines where stages share elements or are correlated through shared
randomness, the theorem requires additional justification or does not apply.

\subsection*{Worked numerical example: a 4-stage AND chain}

To make the composition theorem concrete before we deploy it in Miller-Rabin,
here is a worked example with four stages.

Let $\tilde{b}_1, \tilde{b}_2, \tilde{b}_3, \tilde{b}_4$ be independent
$\Bernoulli{\BoolType}$ instances, each a positive approximate set (FNR $= 0$),
with false-positive rates:
\[
  \fprate_1 = 0.20, \quad \fprate_2 = 0.15, \quad
  \fprate_3 = 0.10, \quad \fprate_4 = 0.25.
\]
We take their four-fold AND: declare $\top$ only if all four report $\top$.

Since $\fnrate_i = 0$ for all stages (positive approximate set), the
four-fold AND has $\fnrate_\wedge = 0$ as well: the AND never falsely
suppresses a true $\top$ (every stage correctly passes $\top$ inputs through).

The composition theorem gives the false-positive rate of the AND. With
$\eta_i = \fprate_i$:
\[
  \tnrate_{\mathrm{total}}
  = (1 - 0.20)(1 - 0.15)(1 - 0.10)(1 - 0.25)
  = 0.80 \times 0.85 \times 0.90 \times 0.75.
\]
Computing step by step:
\begin{align*}
  0.80 \times 0.85 &= 0.680, \\
  0.680 \times 0.90 &= 0.612, \\
  0.612 \times 0.75 &= 0.459.
\end{align*}
So $\tnrate_{\mathrm{total}} = 0.459$ and:
\[
  \fprate_{\wedge,4}
  = 1 - 0.459
  = 0.541.
\]

Is this better than each individual stage? Yes: the individual stages have
$\fprate_i \in \{0.10, 0.15, 0.20, 0.25\}$. The AND gate has
$\fprate_{\wedge,4} \approx 0.541$---actually worse than each individual
stage on its own.

This counterintuitive result is not a paradox. The stages in this example are
quite noisy ($\fprate_i$ up to 0.25). The AND of four noisy stages accumulates
the FNR in the positive direction (here zero, since FNR $= 0$) and pushes the
FPR down \emph{only} when $\fprate_i < 0.5$ and the product
$\prod(1 - \fprate_i)$ is large. With four stages each above $0.10$, the
product of TNRs is only $0.459$, so the AND gate's FPR is $0.541$. To get an
AND gate with FPR below any individual stage, we need the product of individual
TNRs to exceed the best individual TNR, i.e., the other stages must compensate.

For comparison, if all four stages had $\fprate = 0.25$:
\[
  \fprate_{\wedge,4} = 1 - (0.75)^4 = 1 - 0.3164 = 0.684.
\]
And if all four stages had $\fprate = 0.05$:
\[
  \fprate_{\wedge,4} = 1 - (0.95)^4 = 1 - 0.8145 = 0.185.
\]
With low per-stage FPR, the AND gate's FPR is substantially lower than any
individual stage. This is the regime Miller-Rabin exploits: with
$\fprate \leq 1/4$ per round, a $k$-round AND achieves $\fprate_k \leq
(1/4)^k$.

\subsection*{Looking ahead to Miller-Rabin}

The composition theorem establishes the machinery. \Cref{sec:miller-rabin}
deploys it against a canonical example: Miller-Rabin witnesses are
$\Bernoulli{\BoolType}$ instances with FNR $= 0$ and FPR $\leq 1/4$, so
$k$ independent rounds give the bound $\fprate_k \leq (1/4)^k$
from~\eqref{eq:composition-theorem}, a direct instance of the framework
developed in this section.
